
\documentclass[
	% -- opções da classe memoir --
	12pt,				% tamanho da fonte
	openright,			% capítulos começam em pág ímpar (insere página vazia caso preciso)
	oneside,			% para impressão em verso e anverso. Oposto a oneside
	a4paper,			% tamanho do papel. 
	% -- opções da classe abntex2 --
	%chapter=TITLE,		% títulos de capítulos convertidos em letras maiúsculas
	%section=TITLE,		% títulos de seções convertidos em letras maiúsculas
	%subsection=TITLE,	% títulos de subseções convertidos em letras maiúsculas
	%subsubsection=TITLE,% títulos de subsubseções convertidos em letras maiúsculas
	% -- opções do pacote babel --
	english,			% idioma adicional para hifenização
	brazil				% o último idioma é o principal do documento
	]{abntex2}

% ---
% Pacotes básicos 
% ---
\usepackage{unitins}
\usepackage{lmodern}			% Usa a fonte Latin Modern			
\usepackage[T1]{fontenc}		% Selecao de codigos de fonte.
\usepackage[utf8]{inputenc}		% Codificacao do documento (conversão automática dos acentos)
\usepackage{lastpage}			% Usado pela Ficha catalográfica
\usepackage{indentfirst}		% Indenta o primeiro parágrafo de cada seção.
\usepackage{color}				% Controle das cores

\usepackage{graphicx}			% Inclusão de gráficos
% \usepackage{svg}
% \usepackage{xcolor}
\usepackage{subfig}
\usepackage{microtype} 			% para melhorias de justificação
\usepackage{soul}
\usepackage{amssymb}
\usepackage{amsmath}
\usepackage{spreadtab}
\usepackage{multirow}
\usepackage{amsthm}
\usepackage{url}
\usepackage[portuguese, ruled, linesnumbered]{algorithm2e}
\usepackage{float}
\usepackage{appendix}
\usepackage{fancyhdr}
\usepackage{lastpage}
\usepackage{verbatim}

% ---

% ---
% Pacotes de citações
% ---
\usepackage{hyperref}

%\usepackage[brazilian,hyperpageref]{backref}	 % Paginas com as citações na bibl
\usepackage[alf,abnt-etal-list=2 ]{abntex2cite}	% Citações padrão ABNT


\usepackage[table]{xcolor}
\definecolor{lightgray}{gray}{0.9}
\graphicspath{{imagens/}}

\theoremstyle{theorem}
\newtheorem{teo}{Teorema}[chapter]
\newtheorem{lema}[teo]{Lema}
\theoremstyle{definition}
\newtheorem{defi}[teo]{Definição}

\usepackage{svg}
\svgpath{{imagens/svgimgs}} % <- using \svgpath to avoid warning
\usepackage{amsmath}
\usepackage{url}



\usepackage{booktabs}

% ---
% Informações de dados para CAPA e FOLHA DE ROSTO
% ---

\titulo{Governança de TI em Municípios de Baixa Maturidade: proposta de um guia mínimo para elevação do IGovTI nas gestões municipais do Tocantins}
\autor{John Helder Cardoso Alves}
\local{Palmas}
\data{2026}
\orientador{Prof. Me. Tamirys Virgulino Ribeiro Prado}
\instituicao{%
Fundação Universidade do Tocantins
\par
  Curso de Sistemas de Informação
  }
\tipotrabalho{Trabalho de Conclusão de Curso (Graduação)}
% O preambulo deve conter o tipo do trabalho, o objetivo, 
% o nome da instituição e a área de concentração 
\preambulo{
Trabalho apresentado ao Curso de Sistemas de Infor-
mação da Fundação Universidade do Tocantins -
UNITINS como parte dos requisitos para a obtenção
do grau de Bacharel em Sistemas de Informação,
sob a orientação da professora Me. Tamirys Virgulino Ribeiro Prado
}

% ---


% ---
% Configurações de aparência do PDF final

% alterando o aspecto da cor azul
\definecolor{blue}{RGB}{41,5,195}

% informações do PDF
\makeatletter
\hypersetup{
     	%pagebackref=true,
		pdftitle={\@title}, 
		pdfauthor={\@author},
    	pdfsubject={\imprimirpreambulo},
	    pdfcreator={LaTeX with abnTeX2},
		pdfkeywords={Algoritmo}{trabalho acadêmico}, 
		colorlinks=true,       		% false: boxed links; true: colored links
    	linkcolor=blue,          	% color of internal links
    	citecolor=blue,        		% color of links to bibliography
    	filecolor=magenta,      		% color of file links
		urlcolor=blue,
		bookmarksdepth=4
}
\makeatother
% --- 

% --- 
% Espaçamentos entre linhas e parágrafos 
% --- 

% O tamanho do parágrafo é dado por:
\setlength{\parindent}{1.3cm}

% Controle do espaçamento entre um parágrafo e outro:
\setlength{\parskip}{0.2cm}  % tente também \onelineskip

% ---
% compila o indice
% ---
\makeindex
% ---

% ----
% Início do documento
% ----
\begin{document}

% ---
% Definição de comandos pessoais para auxiliar na produção do trabalho
% ---

% COR PERSONALIZADA:
\definecolor{CorTextoNovo}{HTML}{006400}
\definecolor{CorTextoPendente}{HTML}{FFC222}

% \recado{mensagem}
%       - Finalidade: Inserir notas, observações ou dúvidas diretas para o orientador/revisor.
%       - Formatação: Texto pequeno, em negrito e na cor vermelha destacados entre colchetes.
%       - Exemplo: \recado{Conferir se este cálculo está correto.}
\newcommand{\recado}[1]{\small\textcolor{red}{\textbf{[RECADO: #1]}}}


% \revisar{texto}
%       - Finalidade: Destacar blocos de texto recém-escritos que ainda NÃO foram
%                     validados pelo orientador/corretor.
%       - Formatação: Exibe o texto em verde escuro.
%       - Exemplo: \revisar{Esta nova metodologia foi aplicada para resolver o problema X.}
\newcommand{\revisar}[1]{\textcolor{CorTextoNovo}{#1}}


% \revisar{texto}
%       - Finalidade: Destacar blocos de texto recém-escritos que ainda NÃO foram
%                     validados pelo orientador/corretor.
%       - Formatação: Exibe o texto em verde escuro.
%       - Exemplo: \revisar{Esta nova metodologia foi aplicada para resolver o problema X.}
\newcommand{\revisando}[1]{\textcolor{CorTextoPendente}{#1}}


% \revisado{texto}
%       - Finalidade: Marcar trechos que já foram lidos e aprovados.
%       - Formatação: Mantém o texto no formato padrão do documento (sem cor/destaque).
%       - Dica de fluxo: Quando o orientador aprovar o bloco de texto contido em \revisar{...},
%                        basta trocar a palavra "revisar" por "revisado" no código.
\newcommand{\revisado}[1]{#1}


% Retira espaço extra obsoleto entre as frases.
\frenchspacing 

% ----------------------------------------------------------
% ELEMENTOS PRÉ-TEXTUAIS
% ----------------------------------------------------------
% \pretextual
% ----------------------------------------------------------

% % ----------------------------------------------------------
% ELEMENTOS PRÉ-TEXTUAIS
% ----------------------------------------------------------
% \pretextual

% ---
% Capa
% ---
\imprimircapa
% ---

% ---
% Folha de rosto
% (o * indica que haverá a ficha bibliográfica)
% ---
\imprimirfolhaderosto
% ---

% ---
% Inserir folha de aprovação
% ---

% Isto é um exemplo de Folha de aprovação, elemento obrigatório da NBR
% 14724/2011 (seção 4.2.1.3). Você pode utilizar este modelo até a aprovação
% do trabalho. Após isso, substitua todo o conteúdo deste arquivo por uma
% imagem da página assinada pela banca com o comando abaixo:
%
% \includepdf{folhadeaprovacao_final.pdf}
%
\begin{folhadeaprovacao}

  	
  	
  	\begin{center}
  		\includegraphics[width=1\textwidth]{imagens/unitins.png}
  		\ABNTEXchapterfont\Large   CURSO DE SISTEMAS DE INFORMA{\c{C}}{\~{A}}O
  		
  		\par
  		\vspace*{1cm}     
  		{\ABNTEXchapterfont\bfseries\large \expandafter\MakeUppercase{\imprimirtitulo}  \vspace*{1cm}    }
  		\par
  		{\large \expandafter\MakeUppercase{\imprimirautor}}
  		%\vspace*{\fill}
  		\par
  		\vspace*{1cm}     
  		\hspace{.45\textwidth}
  		\begin{minipage}{.5\textwidth}
  			\small\imprimirpreambulo
  			
  		\end{minipage}%
  	%	\vspace*{\fill}
  	\end{center}
 
  
        
  
   \assinatura{\textbf{\imprimirorientador} \\ Orientador} 
   \assinatura{\textbf{Convidado} \\ Convidado 1}
   \assinatura{\textbf{Convidado} \\ Convidado 2}
   %\assinatura{\textbf{Professor} \\ Convidado 3}
   %\assinatura{\textbf{Professor} \\ Convidado 4}
      
   \begin{center}
    \vspace*{0.5cm}
    {\large\imprimirlocal}
    \par
    {\large\imprimirdata}
  \end{center}

% \begin{center}
%     \centering
%     \includegraphics[scale=0.8]{ficha.pdf} 
% \end{center}

% \begin{center}
%     \centering
%     \includegraphics[scale=0.8]{ata_assinada.pdf} 
% \end{center}

\end{folhadeaprovacao}
% ---

% ---
% Dedicatória
% ---
% \begin{dedicatoria}

% \end{dedicatoria}
% ---

% ---
% Agradecimentos
% --- Es
% \begin{agradecimentos}

% \end{agradecimentos}
% ---

% ---
% Epígrafe
% ---
\begin{epigrafe}
    \vspace*{\fill}
	\begin{flushright}
		\textit{``Não vos amoldeis às estruturas deste mundo, \\
		mas transformai-vos pela renovação da mente, \\
		a fim de distinguir qual é a vontade de Deus: \\
		o que é bom, o que Lhe é agradável, o que é perfeito.”\\
		(Bíblia Sagrada, Romanos 12, 2)}
	\end{flushright}
\end{epigrafe}
% ---

% ---
% RESUMOS
% ---

% resumo em português
\setlength{\absparsep}{18pt} % ajusta o espaçamento dos parágrafos do resumo
\begin{resumo}
Lorem Ipsum is simply dummy text of the printing and typesetting industry. Lorem Ipsum has been the industry's standard dummy text ever since 1966, when designers at Letraset and James Mosley, the librarian at St Bride Printing Library in London, took a 1914 Cicero translation and scrambled it to make dummy text for Letraset's Body Type sheets. It has survived not only many decades, but also the leap into electronic typesetting, remaining essentially unchanged. It was popularised thanks to these sheets and more recently with desktop publishing software like Aldus PageMaker and Microsoft Word including versions of Lorem Ipsum.
Lorem Ipsum is simply dummy text of the printing and typesetting industry. Lorem Ipsum has been the industry's standard dummy text ever since 1966, when designers at Letraset and James Mosley, the librarian at St Bride Printing Library in London, took a 1914 Cicero translation and scrambled it to make dummy text for Letraset's Body Type sheets. It has survived not only many decades, but also the leap into electronic typesetting, remaining essentially unchanged. It was popularised thanks to these sheets and more recently with desktop publishing software like Aldus PageMaker and Microsoft Word including versions of Lorem Ipsum.
Lorem Ipsum is simply dummy text of the printing and typesetting industry. Lorem Ipsum has been the industry's standard dummy text ever since 1966, when designers at Letraset and James Mosley, the librarian at St Bride Printing Library in London, took a 1914 Cicero translation and scrambled it to make dummy text for Letraset's Body Type sheets. It has survived not only many decades, but also the leap into electronic typesetting, remaining essentially unchanged. It was popularised thanks to these sheets and more recently with desktop publishing software like Aldus PageMaker and Microsoft Word including versions of Lorem Ipsum.

\textbf{Palavras-chaves}: Lorem, Ipsum, is e simply.
 
\end{resumo}

% resumo em inglês
\begin{resumo}[Abstract]
 \begin{otherlanguage*}{english}
Lorem Ipsum is simply dummy text of the printing and typesetting industry. Lorem Ipsum has been the industry's standard dummy text ever since 1966, when designers at Letraset and James Mosley, the librarian at St Bride Printing Library in London, took a 1914 Cicero translation and scrambled it to make dummy text for Letraset's Body Type sheets. It has survived not only many decades, but also the leap into electronic typesetting, remaining essentially unchanged. It was popularised thanks to these sheets and more recently with desktop publishing software like Aldus PageMaker and Microsoft Word including versions of Lorem Ipsum.
Lorem Ipsum is simply dummy text of the printing and typesetting industry. Lorem Ipsum has been the industry's standard dummy text ever since 1966, when designers at Letraset and James Mosley, the librarian at St Bride Printing Library in London, took a 1914 Cicero translation and scrambled it to make dummy text for Letraset's Body Type sheets. It has survived not only many decades, but also the leap into electronic typesetting, remaining essentially unchanged. It was popularised thanks to these sheets and more recently with desktop publishing software like Aldus PageMaker and Microsoft Word including versions of Lorem Ipsum.
Lorem Ipsum is simply dummy text of the printing and typesetting industry. Lorem Ipsum has been the industry's standard dummy text ever since 1966, when designers at Letraset and James Mosley, the librarian at St Bride Printing Library in London, took a 1914 Cicero translation and scrambled it to make dummy text for Letraset's Body Type sheets. It has survived not only many decades, but also the leap into electronic typesetting, remaining essentially unchanged. It was popularised thanks to these sheets and more recently with desktop publishing software like Aldus PageMaker and Microsoft Word including versions of Lorem Ipsum.



   \vspace{\onelineskip}
 
   \noindent 
   \textbf{Key-words}: Lorem, Ipsum, is e simply.
 \end{otherlanguage*}
\end{resumo}


% ---
% inserir lista de ilustrações
% ---
\pdfbookmark[0]{\listfigurename}{lof}
\listoffigures*
% ---

% ---
% inserir lista de tabelas
% ---
\newpage
\pdfbookmark[0]{\listtablename}{lot}
\newpage
\listoftables*
\cleardoublepage
% ---

% ---
% inserir lista de abreviaturas e siglas
% ---
\begin{siglas}
  
\end{siglas}
% ---


% ---
% inserir o sumario
% ---
\pdfbookmark[0]{\contentsname}{toc}
\tableofcontents*
\cleardoublepage
% ---



% ELEMENTOS TEXTUAIS
% ---------------------------------[!htb]-------------------------
%\textual

%Capitulos

% 
\mainmatter
\chapter{Introdução}\label{introducao}
\section{Objetivos}
\subsection{Objetivo Geral}
O objetivo geral deste trabalho consiste em propor um guia de diretrizes mínimas de governança de TI 
para municípios de baixa maturidade no Estado do Tocantins. Para isso, utiliza-se o i-Gov TI como 
indicador para identificar lacunas prioritárias e subsidiar a elevação dos níveis de governança 
nessas administrações. 

\subsection{Objetivos Específicos}
\begin{itemize}
  \item Classificar critérios de avaliação do I-Gov TI de acordo com os mecanismos de governança de 
  TI.

  \item Mapear as principais deficiências e lacunas de Governança de TI nos municípios com baixos 
  índices no I-Gov TI no estado do Tocantins.

  \item Propor guia de decisões mínimas de governança de TI  com foco na elevação do índice de 
  Governança de TI (I-Gov TI)
\end{itemize}

\section{Justificativa}
A governança pública e a gestão estratégica da Tecnologia da Informação (TI) tornaram-se pilares 
indispensáveis para a eficiência administrativa, a transparência e a prestação de serviços públicos 
de qualidade à sociedade. Contudo, a realidade da gestão municipal frequentemente reflete desafios 
severos relacionados à baixa maturidade de governança, o que compromete o desenvolvimento local e a 
correta aplicação dos recursos públicos.

A relevância e a urgência do tema no âmbito do Estado do Tocantins sustentam-se em dados empíricos 
apurados no ano de 2023 pelo Índice de Efetividade da Gestão Municipal (IEGM). Conforme os dados 
consolidados da pesquisa, que contou com a participação de 123 municípios tocantinenses, o estado 
atingiu uma pontuação média de apenas 33. A escala utilizada pelo IEGM vai de 0 à 100\cite{lima_influencia_2023}. 
Essa pontuação classifica o desempenho geral dos municípios tocantinenses na faixa "C", patamar mais 
baixo na escala de avaliação de efetividade.

O IEGM classifica seus resultados em 5 notas onde cada nota representa uma faixa. A tabela abaixo
demonstra cada nota, a faixa que ela representa e o critério de enquadramento.

\begin{table}[htb]
  \centering
  \caption{Notas de enquadramento do IEGM}
  \label{table:classificacao_iegm}
  \begin{tabular}{clr}
    \toprule
    \textbf{Nota} & \textbf{Faixa} & \textbf{Critério} \\
    \midrule
    A & Altamente Efetivo & Índice superior ou igual a 90\% \\
    B+ & Muito Efetivo & Índice superior ou igual a 75\% e menor que 90\% \\
    B & Efetivo & Índice superior ou igual a 60\% e menor que 75\% \\
    C+ & Em fase de adequação & Índice superior ou igual a 50\% e menor que 60\% \\
    C & Baixo nível de adequação & Índice inferior a 50\% \\
    \bottomrule
  \end{tabular}

  \fonte{Adaptado de \citeonline{silva_alise_2024}.}
\end{table}

A partir do levantamento e da análise individual dos 123 municípios participantes realizados nesta 
pesquisa com base nas informações disponibilizadas pelo portal do IEGM (2023), identificou-se que a 
gravidade do cenário revela-se ainda mais proeminente quando se analisa o enquadramento individual 
dos municípios:   

\begin{itemize}
  \item Faixa C (Baixa Efetividade): Absoluta maioria do estado, englobando 113 municípios 
  (91,86\% do total analisado), dos quais apenas 7 alcançaram o nível C+.

  \item Faixa B (Efetividade Média): Compreende apenas 8 municípios, sendo que somente 1 obteve a 
  qualificação B+.

  \item Faixa A (Efetividade Elevada): Restrita a apenas 2 municípios
\end{itemize}

Estes dados evidenciam uma vulnerabilidade estrutural generalizada. A concentração expressiva de 
mais de 90\% das gestões locais na faixa inferior (C) comprova que os municípios do Tocantins 
enfrentam limitações graves no estabelecimento de mecanismos consistentes de governança.

Diante deste diagnóstico, a realização deste trabalho justifica-se pela necessidade crítica de 
compreender as causas dessa baixa maturidade e propor diretrizes mínimas aplicáveis à realidade das 
administrações municipais. Ao propor um guia norteador adequado à realidade local, a pesquisa 
contribui para o fortalecimento institucional, a gestão mais eficiente dos recursos públicos e a 
melhoria dos serviços prestados à população tocantinense. Isto, por consequência, resulta no avanço 
dos indicadores de efetividade e governança de TI do estado.
\chapter{Referencial Teórico}
O objetivo desta seção é fundamentar teoricamente a pesquisa, delimitando os conceitos, princípios e 
parâmetros institucionais que norteiam o estudo. Este capítulo visa conferir rigor metodológico e 
consistência analítica à avaliação da gestão pública municipal e às suas ferramentas de governança.

\section{Governança de TI}

\recado{Citar fonte mais diretamente}
\revisando{
A governança de TI pode ser compreendida como o conjunto de estruturas, processos e mecanismos 
institucionais estruturados para definir direitos decisórios, responsabilidades e formas de 
controle, visando garantir o pleno alinhamento do uso da tecnologia aos objetivos organizacionais \cite{governanca_ia_2026}.
}

\revisado{
Nessa perspectiva, ela estabelece a especificação dos direitos de decisão e o arcabouço de 
responsabilidades necessários para incentivar comportamentos desejáveis no uso da tecnologia. Assim, 
a governança de TI reflete os princípios amplos da governança corporativa, focando na administração 
e no direcionamento da TI para o alcance das metas de desempenho da organização \cite{parcial_2025_processo_gov_ti}.
}

\revisado{
Sob a ótica da alta administração, o termo descreve a forma como os líderes e responsáveis pela 
governança organizacional consideram a TI nas funções de supervisão, monitoramento, controle e 
direcionamento. A maneira como a tecnologia é aplicada impacta diretamente a capacidade da 
organização de atingir sua visão, missão e objetivos estratégicos \cite{rondonia_2023}.
}

\revisado{
Por fim, é essencial diferenciar esse conceito do nível operacional: enquanto a gestão de TI foca na 
execução e na administração cotidiana dos recursos tecnológicos, a governança de TI atua no plano 
estratégico, sendo responsável por definir prioridades, pactuar responsabilidades e instituir 
mecanismos de controle e monitoramento das decisões de tecnologia \cite{governanca_ia_2026}.
}

\revisado{
A governança de TI originou-se da governança corporativa e é composta por estruturas, processos 
organizacionais e elementos de liderança que asseguram a sustentação e a expansão das estratégias e 
metas organizacionais por meio da tecnologia. Por essa razão, prevalece na literatura o entendimento 
de que a governança de TI constitui um subconjunto da governança corporativa \cite{rangel_2021}.
}

\revisado{
O fortalecimento dessa relação foi impulsionado pelo surgimento de normativas rígidas, a exemplo do 
Ato Sarbanes-Oxley, promulgado nos Estados Unidos em 2002, que forçou diversas organizações a 
reestruturarem suas práticas de governança corporativa para assegurar controles fiscais adequados 
aos acionistas e demais partes interessadas \cite{rangel_2021}.
}

\revisado{
Nesse contexto, a governança viabiliza o alinhamento entre a TI e o negócio, garantindo que as 
deliberações da alta direção considerem aspectos como gestão de riscos, planejamento financeiro, 
continuidade dos negócios e conformidade regulatória. Com isso, a TI deixa de ser percebida como um 
mero centro de custos ou instrumento de suporte e passa a integrar a gestão estratégica e financeira 
da organização \cite{joyce_mariana_pereira_2023_7555947}.
}

\revisado{
A governança é responsável por assegurar que as necessidades, condições e opções das partes 
interessadas sejam devidamente avaliadas para estabelecer objetivos corporativos equilibrados e 
pactuados. A partir disso, define-se o direcionamento estratégico mediante priorizações e tomadas de 
decisão, com o subsequente monitoramento do desempenho e da conformidade frente aos alvos fixados 
\cite{joyce_mariana_pereira_2023_7555947}.
}

\revisado{
Para atingir esses objetivos, a Governança de TI estrutura-se em áreas de foco interconectadas, 
como alinhamento estratégico, entrega de valor, gestão de riscos, gestão de recursos e mensuração de 
desempenho, sendo a aplicação conjunta dessas etapas indispensável para a obtenção dos melhores 
resultados organizacionais \cite{rondonia_2023}.
}

\revisado{
Nesse contexto, destacam-se dimensões fundamentais para a consolidação do papel da tecnologia no 
negócio. O alinhamento estratégico assegura que as ações de TI promovam a missão organizacional, 
enquanto a entrega de valor garante a execução de projetos dentro do prazo e do orçamento previstos. 
Paralelamente, a gestão de recursos otimiza os investimentos e os ativos de informação, a gestão de 
riscos protege o patrimônio tecnológico e informacional da instituição, e a mensuração de desempenho 
monitora continuamente a implementação das estratégias para maximizar o valor gerado pela TI \cite{rangel_2021}.
}

\revisado{
Para que os objetivos e áreas de foco da tecnologia sejam efetivamente operacionalizados, a 
literatura consolida que os mecanismos de governança de TI se organizam em três dimensões 
complementares: estruturas, processos e relacionamentos. As estruturas englobam os comitês de 
governança digital, unidades de TI, escritórios de projetos, áreas de segurança e fóruns decisórios. 
Os processos referem-se ao planejamento, priorização de demandas, gestão de portfólio, riscos e 
benefícios, além do acompanhamento de indicadores e auditoria. Por fim, os mecanismos relacionais 
conectam a TI às áreas finalísticas por meio de comunicação integrada, capacitação contínua, 
engajamento da alta administração e construção de uma linguagem comum na organização 
\cite{mello_ramos_2026_21402664}.
}

\subsection{Governança de TI na Administração Pública}

\revisado{
A aplicação da governança no contexto governamental exige compreender que as estratégias do setor 
público diferem substancialmente das do setor privado. Enquanto as empresas privadas buscam o 
retorno sobre o capital e o valor aos acionistas, as instituições públicas voltam-se à entrega de 
valor público. Dessa forma, os investimentos estatais em tecnologia devem ser avaliados com base na 
efetividade do atendimento à sociedade e no cumprimento das metas fixadas nos programas 
governamentais, extrapolando métricas puramente econômicas \cite{neto_avaliacao_2020}.
}

\revisado{
Embora compartilhem pilares conceituais básicos, os modelos de governança de TI em ambos os setores 
distinguem-se em razão do ambiente institucional. No setor público, prevalecem exigências de 
transparência, legitimidade das ações perante a sociedade e conformidade a dispositivos legais e 
órgãos de controle, que sobrepõem-se ao foco em resultados financeiros característico da iniciativa 
privada \cite{sena_governancti_2018}.
}

\revisado{
Essas particularidades ambientais manifestam-se em procedimentos burocráticos rígidos, tomadas de 
decisão altamente estruturadas, dependência de políticas públicas, elevada responsabilização da 
administração e rotatividade de lideranças em virtude de cargos temporários. Esse conjunto de 
contingências afeta diretamente a condução da tecnologia e demanda arranjos de governança 
ajustados à cultura do setor público \cite{sena_governancti_2018}.
}

\revisado{
Nesse cenário, a condução da governança de TI constitui responsabilidade indelegável da Alta 
Administração - compreendendo a liderança política e o corpo executivo. Cabe aos ocupantes do topo 
da organização fixar diretrizes decisórias, gerir riscos e coordenar estrategicamente as ações 
tecnológicas, evitando a falha recorrente de transferir o papel de governança estritamente para a 
área operacional de tecnologia \cite{sena_governancti_2018}.
}

\revisado{
Entretanto, as instituições públicas enfrentam entraves para efetivar essa gestão. Na administração 
pública, a escassez de recursos humanos e de profissionais com competências gerenciais dificulta a 
separação entre as atribuições estratégicas de governança e as rotinas operacionais de gestão. 
Além disso, o foco predominantemente técnico dos concursos públicos resulta na escassez de 
habilidades de planejamento e governança entre os servidores \cite{torriani_modelo_2019}.
}

\recado{Citar fonte mais diretamente}
\revisando{
Outro desafio reside no fato de que frameworks e padrões de boas práticas amplamente difundidos no 
mercado, como COBIT e ITIL, foram concebidos sob a ótica de organizações comerciais lucrativas, 
colidindo com as restrições orçamentárias e a atividade-fim governamental. Assim, a governança 
pública requer a simplificação e a adaptação dessas estruturas à realidade estatal para garantir a 
viabilidade das práticas \cite{sengik_governanctecnologia_2022}.
}

\revisado{
Em razão dessas barreiras, a presença de modelos formais e frameworks internacionais não assegura, 
por si só, uma governança efetiva no setor público. É comum observar a implementação parcial ou 
meramente documental de mecanismos de TI. Os obstáculos consolidados na literatura incluem escassez 
de recursos, baixo engajamento da alta administração, limitações na gestão e reduzida maturidade 
organizacional \cite{sena_governancti_2018}.
}

% { Removido -
% Para induzir a superação desses déficits na administração pública brasileira, o Tribunal de Contas 
% da União (TCU) instituiu levantamentos periódicos de avaliação e monitoramento. Essa medição 
% sistematizada de indicadores de desempenho visa mensurar a maturidade da governança de TI e orientar 
% as instituições públicas na busca por resultados e no aprimoramento da prestação de contas à 
% sociedade \cite{neto_avaliacao_2020}.
% }

%% colocar uma fala de conexão dos paragrafos

\revisar{
A avaliação e o monitoramento da governança pública pelo Tribunal de Contas da União tiveram início 
com levantamentos de perfil focados em Tecnologia da Informação, delineados com o propósito 
explícito de apoiar a melhoria contínua do desempenho das organizações públicas \cite{tcu_referencial_2020}. 
}

\revisar{
Realizados por meio de diagnósticos e questionários eletrônicos, esses levantamentos geram 
indicadores que permitem às instituições identificar vulnerabilidades em suas estruturas, subsidiar 
a definição de objetivos de aprimoramento e priorizar práticas de governança capazes de mitigar 
riscos que comprometam as entregas e a geração de valor para a sociedade \cite{tcu_referencial_2020}.
}

\revisado{
Estudos empíricos aplicados ao setor público nacional corroboram que as dimensões "Estratégias e 
Planos", "Liderança" e "Pessoas" são as variáveis de maior peso na diferenciação dos níveis de 
maturidade de governança de TI nas instituições. Fica evidente, portanto, que os órgãos públicos 
precisam focar esforços prioritários no alinhamento estratégico, na formação contínua e no 
desenvolvimento de competências de seus quadros para elevar seu grau de maturidade 
\cite{campelo_medeiros_maturidade_2016}.
}


\subsection{Mecanismos de Governança de TI}

\revisado{
A operacionalização da governança de TI nas organizações exige a implementação coordenada de um 
conjunto de mecanismos práticos, classificados pela literatura em três dimensões complementares: 
estruturais, processuais e relacionais. Enquanto as deliberações de alto nível fixam diretrizes e 
prioridades, cabe a esses mecanismos operacionalizar a execução cotidiana e o controle das decisões 
tecnológicas, permitindo a efetiva transição do plano estratégico para a prática administrativa 
\cite{governanca_ia_2026}.
}

\revisando{
No modelo de governança de TI proposto por \cite{sengik_governanctecnologia_2022} para a 
Administração Pública, esses mecanismos são organizados no Bloco de Operacionalização em três eixos 
de ação institucional.
}


\begin{itemize}
    \item \revisando{\textbf{Mecanismos Estruturais}: referem-se aos arranjos organizacionais formais e comitês 
    institucionais encarregados de conectar a alta administração, as áreas finalísticas e a área de 
    TI. Incluem a criação do Comitê Diretivo de TI, a especificação clara de papéis e 
    responsabilidades decisórias e a participação formal dos gestores de TI em instâncias superiores 
    de decisão.}

    \item \revisando{\textbf{Mecanismos Processuais}: envolvem a formalização, padronização e institucionalização das 
    rotinas de tomada de decisão e monitoramento das ações tecnológicas. Abrangem a definição de 
    indicadores de desempenho, o Planejamento Estratégico de TI, a formulação de Políticas de 
    Segurança da Informação, as práticas de gestão de riscos, a gestão do portfólio de projetos, a 
    priorização de investimentos e a gestão de serviços de TI.}
    
    \item \revisando{\textbf{Mecanismos Relacionais}: englobam os instrumentos formais e informais direcionados a 
    promover a comunicação, a liderança, o aprendizado compartilhado e a colaboração ativa entre os 
    atores de TI e os gestores de negócio. Contemplam a colaboração entre stakeholders, o 
    compartilhamento de conhecimento, a comunicação efetiva, o exercício da liderança de TI e as 
    parcerias estratégicas entre a TI e as áreas finalísticas.}
\end{itemize}


\revisado{
A eficácia do modelo depende da articulação harmônica entre esses três tipos de mecanismos e as 
áreas de foco da governança de TI, a exemplo do alinhamento estratégico, da gestão de riscos, da 
entrega de valor, do gerenciamento de recursos, da mensuração do desempenho e da prestação de contas. 
Nenhum mecanismo funciona de modo isolado: o desempenho de uma instância estrutural está 
condicionado ao suporte de processos bem delimitados e à fluidez dos canais relacionais que garantem 
o engajamento das partes interessadas \cite{sengik_governanctecnologia_2022}.
}

\revisado{
Com isso, a adoção da taxonomia de mecanismos estruturais, processuais e relacionais estabelecida 
por \cite{sengik_governanctecnologia_2022} fornece um arcabouço analítico sistemático para diagnosticar e avaliar a 
gestão pública. Essa abordagem possibilita verificar não apenas a existência de instâncias formais, 
mas também a maturidade das rotinas operacionais e a efetividade do relacionamento institucional 
indispensável para a geração de valor público por meio da tecnologia 
\cite{sengik_governanctecnologia_2022}.
}

\section{Índice de Efetividade da Gestão Municipal}
Concebido originalmente em 2014 pelo Tribunal de Contas do Estado de São Paulo (TCE-SP), o Índice de 
Efetividade da Gestão Municipal (IEGM) foi posteriormente incorporado e aperfeiçoado pelo Instituto 
Rui Barbosa (IRB) - entidade voltada ao desenvolvimento de estudos e pesquisas para o sistema de 
controle externo \cite{iegm_ceara_2022}. Com o intuito de uniformizar a avaliação das administrações 
públicas locais, o índice foi selecionado em 2016 pela Rede Nacional de Indicadores Públicos 
(Indicon) para aplicação em âmbito nacional \cite{iegm_ceara_2022}.

O indicador foi projetado para mensurar o nível de eficácia e desenvolvimento das 
administrações municipais brasileiras. O IEGM engloba sete áreas fundamentais: educação, saúde, 
gestão fiscal, planejamento, meio ambiente, defesa civil e governança de TI \cite{lima_influencia_2023}.

O resultado geral do IEGM varia de 0 a 100, no qual a pontuação máxima representa a aderência 
plena do município às normas e aos controles estabelecidos. Para o cálculo do índice consolidado, 
aplica-se uma métrica de ponderação que reflete a relevância de cada dimensão específica. Assim, a 
composição do IEGM é estruturada pelas seguintes proporções:
\cite{lima_influencia_2023}.


\begin{itemize}
  \item Educação (i-Educ, 20\%)
  \item Saúde (i-Saúde, 20\%)
  \item Gestão Fiscal (i-Fiscal, 20\%)
  \item \revisar{Planejamento (i-Planejamento, 20\%)}
  \item Meio Ambiente (i-Amb, 10\%)
  \item Proteção dos Cidadãos/Defesa Civil (i-Cidade, 5\%)
  \item Governança de TI (i-Gov TI, 5\%)
\end{itemize}

O Índice conta com a adesão de todos os 27 Tribunais de Contas dos estados e do Distrito Federal, 
consolidando-se como um instrumento nacional de avaliação das administrações públicas locais \cite{araujo_gestoes_municipais_2021}. 
A operacionalização do índice ocorre por meio do Termo de Adesão à Rede de Indicadores Nacionais, 
firmado entre o Instituto Rui Barbosa (IRB) e os órgãos de contas estaduais e distrital. A partir 
desse acordo, as informações prestadas pelos municípios são coletadas e posteriormente validadas 
in loco ou auditadas pelos respectivos tribunais, garantindo a consistência dos dados do índice \cite{araujo_gestoes_municipais_2021}.

\subsection{Índice de Governança de TI (I-Gov TI)}

\revisar{No contexto da transformação digital do setor público, a governança de Tecnologia da Informação 
posiciona-se como elemento indispensável para a otimização dos serviços estatais. Ao impulsionar a 
transparência, a integridade e a eficiência das entregas governamentais, o uso estratégico da TI 
reflete diretamente na efetividade da gestão pública. O I-Gov TI surge, portanto, como um indicador 
estruturado para diagnosticar e induzir o aprimoramento da governança tecnológica municipal, em 
estrita observância aos princípios da moralidade e da eficiência administrativa \cite{silva2025implementaccao}.}

\revisar{Como componente estruturante do IEGM, o i-Gov TI mensura o alinhamento entre a gestão de TI e as 
demandas coletivas. Sua avaliação incide sobre os pilares da governança digital municipal: 
diretrizes para uso de TIC, protocolos de segurança da informação, desenvolvimento de competências 
da equipe interna e fomento à transparência governamental \cite{construcao_iegm_2018}.}


\revisar{A aferição desses aspectos reflete de forma direta a capacidade dos municípios em atender às 
exigências normativas e em incorporar as melhores práticas de gestão pública. Cumpre destacar que 
toda essa estrutura avaliativa fundamenta-se rigorosamente nos diplomas legislativos que orientam, 
normatizam e direcionam a boa governança de Tecnologia da Informação \cite{CEPCP_2025}.}

\revisar{Nessa perspectiva, o i-Gov TI alinha-se a tais exigências ao verificar a participação formal da 
equipe de Tecnologia da Informação nos processos licitatórios do setor. A exigência do Plano Diretor
de Tecnologia da Informação (PDTI) para a elaboração de um Estudo Técnico Preliminar (ETP) 
consistente posiciona o planejamento como instrumento central de economicidade e conformidade 
contratual. Corroborando essa diretriz, a jurisprudência do Tribunal de Contas da União (TCU) aponta 
recorrentemente a deficiência no planejamento como causa primária de contratações ineficazes. A 
ausência de planejamento estratégico em TI, portanto, compromete a eficiência dos investimentos 
públicos e expõe os municípios a riscos normativos e à responsabilização de seus gestores \cite{CEPCP_2025}.}

\revisar{Além do respaldo normativo e jurisprudencial, a avaliação do i-Gov TI apoia-se em frameworks 
internacionais de governança de TI. O COBIT, por exemplo, estrutura práticas e processos voltados ao 
alinhamento estratégico da tecnologia com os objetivos institucionais e a geração de valor público. 
A dimensão i-Gov TI incorpora essa perspectiva ao averiguar se os municípios possuem planejamento 
estratégico de TI integrado ao planejamento governamental global, instigando a superação de 
iniciativas isoladas \cite{CEPCP_2025}.}

% \section{Governança de TI na Administração Pública}


% \section{Governança Frugal de TI em Municípios de Baixa Maturidade}


% \section{Design Science Research(DSR) como Abordagem Metodológica}


% \section{Trabalhos Relacionados}




% 
\mainmatter
\chapter{Metodologia da Pesquisa}\label{metodologia}

\section{Descrição dos Experimentos}
Propor um guia de diretrizes mínimas de governança de TI para municípios de baixa maturidade no 
Estado do Tocantins, utilizando o i-GovTI como instrumento diagnóstico. 

\subsection{Experimento 01}
% Classificar critérios de avaliação do 
% I-Gov TI de acordo com os mecanismos de governança de TI 
% definidos no referencial teórico

Categorizar as questões do formulário conforme os mecanismos de governança de TI apresentado no
trabalho realizado por \cite{sengik_governanctecnologia_2022} descrito com três
categorias: Mecanismo estrutural, mecanismos processuais e mecanismos relacionais.


\subsection{Experimento 02}
% \textit{*nota* Mapear as principais deficiências e lacunas de Governança de TI nos municípios com 
% baixos índices no iGovTI no estado do Tocantins"}

Será realizado uma pesquisa quantitativo-descritiva dos dados disponibilizados pelo site do 
Indicador IEGM na forma de questionário. 

A base de dados que será utilizado para diagnóstico foi constituída pelas respostas declaratórias 
dos 123 municípios do Estado do Tocantins ao questionário oficial do Índice de Efetividade da Gestão 
Municipal (IEGM), exclusivamente na dimensão i-GovTI composto por 37 questões.

Os procedimentos de tratamento e avaliação dos dados seguiram três etapas encadeadas:

Recorte Amostral: Não será aplicado nenhum filtro eliminatório, portanto será enquadrado 
todos os municípios e faixas de classificação do I-GovTI.

Análise de Frequência e Ranking de Inadequação: Levantamento do percentual de não conformidade para 
cada quesito, (permitindo ranquear as lacunas mais recorrentes e mapear os gargalos prioritários de 
gestão que serviram como requisitos para a elaboração do Guia de Governança {pensar melhor na
descricao dos resultados esperados do experimento}).

\subsection{Experimento 03}
% \textit{*nota* Propor guia de decisões mínimas de governança de TI  com foco na elevação do índice 
% de Governança de TI(IGovTI)"}
% 
\mainmatter
\chapter{Resultados e Discussão}\label{resultados}




% 
\mainmatter
\chapter{Considerações Finais}\label{conclusao}






% ----------------------------------------------------------
% ELEMENTOS PÓS-TEXTUAIS
% ----------------------------------------------------------
\postextual
% ----------------------------------------------------------

% ----------------------------------------------------------
% Referências bibliográficas
% ----------------------------------------------------------
\bibliography{bibliografia}

% \include{apendices}
% ----------------------------------------------------------
% Glossário
% ----------------------------------------------------------
%
% Consulte o manual da classe abntex2 para orientações sobre o glossário.
%
%\glossary

%---------------------------------------------------------------------
% INDICE REMISSIVO
%---------------------------------------------------------------------
\phantompart
\printindex
%---------------------------------------------------------------------

\end{document}
\grid
